\section{AI 使用声明}\label{sec:ai-decl}

本节说明生成式 AI 在本项目中的使用范围、边界与责任归属，依据是赛题允许使用 AI
但须体现参赛者的设计判断，以及评分标准对 AI 使用声明严谨程度的考察要求。

\subsection{使用范围}

生成式 AI 在本项目中承担辅助性工作：资料检索线索整理、建模方案讨论、代码草拟与调试、
测试用例设计、文字整理、可视化方案构思，以及装饰性视觉素材的生成。

\subsection{边界}

\begin{itemize}
  \item \textbf{不生成实验数据}：本文的实验数据、模型指标、统计图与网络结构图
        均由实际程序运行产生。AI 生成的装饰性图像不作为实验数据或科学证据。
  \item \textbf{不填未定数值}：上游文档未定义的数值、公式与语义，实现侧不得自行选定并当作既定事实。
        本文引用的每个数字均可回溯到运行产物或配置；未跑完的实验留空并注明，不以估算替代。
  \item \textbf{不替代核心设计}：问题定义、数学模型（发育映射、网络动力学、遗传边界与当代学习）、
        实验设计与结论判断均由团队成员完成，AI 仅作为实现与整理的加速工具。
\end{itemize}

\subsection{责任归属}

团队成员对问题定义、数学模型、代码、数据、实验、统计分析、文献引用、
技术报告与最终路演陈述进行人工审核，并承担最终责任。
